pdfoutput=1
\pdfoutput=1
\documentclass[10pt,a4paper]{article}

% --- Tight layout for 2-page fit ---
\usepackage[margin=1.8cm]{geometry}
\usepackage{amsmath,amssymb,amsthm}
\usepackage{enumitem}
\usepackage{hyperref}
\usepackage[numbers,sort&compress]{natbib}

% --- Theorem environments ---
\newtheorem{theorem}{Theorem}[section]
\newtheorem{corollary}[theorem]{Corollary}
\theoremstyle{definition}
\newtheorem{definition}[theorem]{Definition}
\theoremstyle{remark}
\newtheorem{remark}[theorem]{Remark}

% --- Compact spacing ---
\setlist{nosep,leftmargin=*}
\AtBeginDocument{\setlength{\abovedisplayskip}{4pt plus 1pt minus 2pt}}
\AtBeginDocument{\setlength{\belowdisplayskip}{4pt plus 1pt minus 2pt}}
\AtBeginDocument{\setlength{\abovedisplayshortskip}{2pt plus 1pt}}
\AtBeginDocument{\setlength{\belowdisplayshortskip}{2pt plus 1pt}}

\title{\bfseries\large Noise-Difficulty Indistinguishability:\\An Information-Theoretic Lower Bound\\for Data Quality and LLM Hallucination}

\author{SCX}}
\date{2026}

\begin{document}
\maketitle
\thispagestyle{empty}
\vspace{-8mm}

% ===================================================================
\section{Theorem Statement}
% ===================================================================

Let $\mathcal{X}\times\mathcal{Y}$ be the input-label space with $|\mathcal{Y}|=K\geq2$,
and $\{f_m\}_{m=1}^M$ be $M$ experts with conditionally independent predictions given~$x$.
The ground-truth oracle $f^*:\mathcal{X}\to\mathcal{Y}$ is unobservable.

\begin{definition}
Two worlds are considered:
\begin{enumerate}[label=(\roman*)]
  \item $\mathcal{P}_{\text{noise}}$ (World~A): a fraction~$\rho$ of inputs suffers
        \emph{label noise}---the observed label~$y$ differs from $y^*$ with probability~$\eta_{\text{err}}$.
  \item $\mathcal{P}_{\text{hard}}$ (World~B): all $y=y^*$, but the same fraction~$\rho$
        is \emph{intrinsically difficult}---experts err with effective rate~$\eta_{\text{amb}}$.
\end{enumerate}
\end{definition}

\begin{theorem}[Noise-Difficulty Indistinguishability]\label{thm:main}
For any $K\geq2$, $M\geq1$, there exist $\mathcal{P}_{\text{noise}}$ and $\mathcal{P}_{\text{hard}}$ such that:
\begin{enumerate}[label=(\arabic*)]
  \item \textbf{Observational equivalence.}
        $\mathcal{P}_{\text{noise}}(x,y,\{f_m\}) = \mathcal{P}_{\text{hard}}(x,y,\{f_m\})$,
        \emph{iff} $\eta_{\text{err}}=\eta_{\text{amb}}=\eta$.
        \label{item:equiv}

  \item \textbf{Algorithmic lower bound.}
        For any algorithm $\mathcal{A}$ using only observations,
        \begin{equation}\label{eq:bound}
        \max\bigl(\mathrm{Error}_{\mathcal{P}_{\text{noise}}}(\mathcal{A}),\;
                   \mathrm{Error}_{\mathcal{P}_{\text{hard}}}(\mathcal{A})\bigr)
        \ge \frac{\eta\rho}{2}>0.
        \end{equation}
        \label{item:bound}
\end{enumerate}
\end{theorem}

No algorithm can distinguish ``label is wrong'' from ``sample is hard'' using observed
expert outputs alone---the two worlds are observationally identical, and any discrimination
incurs error $\ge\eta\rho/2$.

% ===================================================================
\section{Constructive Proof ($K=2$)}
% ===================================================================

Fix states $s_1,s_2$ with $\mathbb{P}(X\!\in\!s_1)=\rho$, $\mathbb{P}(X\!\in\!s_2)=1-\rho$.
Choose $\eta\!\in\!(0,\tfrac12)$, $\varepsilon_1\!<\!\varepsilon_2\!\in\!(0,\tfrac12)$.

\subsection*{World~A (Noise-driven)}

State $s_1$: $y^*\equiv0$.  Labels flipped with probability $\eta_{\text{err}}$:
$\mathbb{P}(y\!\neq\!y^*\mid s_1)=\eta_{\text{err}}$.
Expert accuracy on clean labels: $\mathbb{P}(f_m=y^*\mid s_1,\text{clean})=1-\varepsilon_1$.
State $s_2$: $y^*\equiv0$, no noise, accuracy $1-\varepsilon_2$.

\subsection*{World~B (Difficulty-driven)}

State $s_1$: $y=y^*$ but $y^*$ is random:
$\mathbb{P}(y^*=0\mid s_1)=1-\eta_{\text{amb}}$,
$\mathbb{P}(y^*=1\mid s_1)=\eta_{\text{amb}}$.
Experts are biased toward~0: $\mathbb{P}(f_m=0\mid s_1,y^*=0)=\mathbb{P}(f_m=0\mid s_1,y^*=1)=1-\varepsilon_1$.
Thus accuracy is $1-\varepsilon_1$ when $y^*=0$ but drops to $\varepsilon_1$ when $y^*=1$.
State $s_2$: identical to World~A.

\subsection*{Equivalence verification}

Under $\eta_{\text{err}}=\eta_{\text{amb}}=\eta$, the observable distribution factorises as
$\mathcal{P}(x,y,\{f_m\})=\mathcal{P}(x)\mathcal{P}(y\mid x)\prod_m\mathcal{P}(f_m\mid x)$:

\smallskip
\textit{Labels:}
$\mathcal{P}_{\text{noise}}(y=0\mid s_1)=(1-\eta)\cdot1+\eta\cdot0=1-\eta
 =\mathcal{P}_{\text{hard}}(y=0\mid s_1)$.

\textit{Experts:}
$\mathcal{P}_{\text{noise}}(f_m=0\mid s_1)=1-\varepsilon_1
 =\mathcal{P}_{\text{hard}}(f_m=0\mid s_1)$
(since $(1-\eta)(1-\varepsilon_1)+\eta(1-\varepsilon_1)=1-\varepsilon_1$).

All marginals match; conditional independence yields identical joint distributions,
proving item~(\ref{item:equiv}).

\subsection*{Algorithmic impossibility}

Let $\mathcal{A}$ flag samples as noisy.  Let $a$ be the fraction of $\{x\!\in\!s_1,y=1\}$
flagged.  In World~A these $\rho\eta$ samples are truly noisy:
$\mathrm{Error}_{\text{noise}}=\rho\eta(1-a)$.
In World~B they are correct but difficult: $\mathrm{Error}_{\text{hard}}=\rho\eta a$ (false alarms).
Hence
$\max(\mathrm{Error}_{\text{noise}},\mathrm{Error}_{\text{hard}})
 \ge\frac{\rho\eta(1-a)+\rho\eta a}{2}=\frac{\rho\eta}{2}$,
proving~(\ref{item:bound}).  Minimax optimum: $a=\tfrac12$ (coin flip). \hfill$\square$

\begin{remark}[$K>2$]
For $K>2$, construct World~B with fully random experts independent of $y^*$ whose
marginals match World~A's.  This extreme construction---experts as random guessers---suffices
as an existence proof. \hfill$\square$
\end{remark}

% ===================================================================
\section{The Symmetry Condition $\eta_{\text{err}}=\eta_{\text{amb}}$}
% ===================================================================

The condition $\eta_{\text{err}}=\eta_{\text{amb}}$ is \emph{necessary and sufficient}
for observational equivalence:
\[
\mathcal{P}_{\text{noise}}(y=1\mid s_1)=\eta_{\text{err}},\qquad
\mathcal{P}_{\text{hard}}(y=1\mid s_1)=\eta_{\text{amb}}.
\]
When $\eta_{\text{err}}\neq\eta_{\text{amb}}$, the two worlds produce distinct label-class
frequencies and become statistically distinguishable---the impossibility collapses.
The symmetry thus delineates \emph{precisely when} noise and difficulty are confounded.
Interpretation: the probability that flipping noise corrupts a label must equal the
probability that ambiguity alone generates an alternative valid label.  This is a
testable condition: compare label marginals across candidate partitions of~$\mathcal{X}$.

% ===================================================================
\section{The Lower Bound $\eta\rho/2$}
% ===================================================================

The bound admits three readings:
\begin{enumerate}[label=(\alph*)]
  \item \textbf{Information-theoretic.} The two worlds differ on $\rho\eta\cdot n$ samples.
        Identical observation distributions force coin-flip decisions on these, giving $\rho\eta/2$.
  \item \textbf{Tightness.} Random labelling ($a=\tfrac12$) achieves exactly $\rho\eta/2$,
        so the bound is minimax-optimal.
  \item \textbf{Positivity.} Strictly positive for any $\eta>0,\rho>0$; perfect
        discrimination requires either zero ambiguity or empty ambiguous subset.
\end{enumerate}
This is a \emph{fundamental} limit---not an artifact of model class or algorithm design.

% ===================================================================
\section{Application: LLM Hallucination}
% ===================================================================

\subsection*{Multi-seed independent decoding}

Run an LLM $M$ times on context~$c$ with independent random seeds $r_m$.  For fixed
model~$\theta$, the deterministic forward pass $f_\theta(c)$ is identical; stochastic
perturbations $\xi_m=\xi(c;r_m)$ (dropout, sampling, temperature) satisfy
$\xi_m\perp\xi_{m'}\!\mid\!c$.  Error indicators $e_m(c)=\mathbf{1}\{\text{output of seed }m\text{ is wrong}\}$
are conditionally independent: $e_m\perp e_{m'}\!\mid\!c$.  The $M$ seeds thus constitute
$M$ conditionally independent experts, and Theorem~\ref{thm:main} applies.

\begin{corollary}[Hallucination lower bound]\label{cor:hall}
For an LLM with $M$ independent seeds,
$\mathrm{Hallucination}_{\text{ambiguous}} \ge \eta\rho/2$,
where $\eta$ is model error probability on ambiguous queries and $\rho$ is their proportion.
\end{corollary}

From seed outputs alone, one cannot distinguish model error from essential query ambiguity.

\subsection*{Symmetry in the LLM setting}

$\eta_{\text{err}}=\eta_{\text{amb}}$ means: the perturbation $\xi_m$ has equal probability
of pushing output away from the correct answer (clear query) as pushing toward a valid
alternative (ambiguous query).  This depends on logit geometry and is empirically testable.

\subsection*{Breaking the bound}

The premise is ``only the $M$ seed outputs.''  External information breaks it:
\textbf{RAG} introduces retrieved documents; \textbf{RLHF} injects human preference anchors;
\textbf{tool use} provides execution verification.
\textbf{Hallucinations cannot be eliminated by scaling alone---only by breaking
the information closure of autoregressive generation.}

% ===================================================================
\section*{Acknowledgments}

Supported by .  Theorem~\ref{thm:main} emerged from the SCX
axiomatic framework \cite{scx2026} but is fully self-contained: it requires only conditional
independence of expert errors.

\begin{thebibliography}{10}

\bibitem{scx2026}
SCX.
\newblock A Taxonomic Theory of Neural Networks.
\newblock arXiv preprint, 2026.

\bibitem{wolpert1996}
D.H.\ Wolpert.
\newblock The Lack of A~Priori Distinctions Between Learning Algorithms.
\newblock \textit{Neural Computation}, 8(7):1341--1390, 1996.

\bibitem{dawid1979}
A.P.\ Dawid and A.M.\ Skene.
\newblock Maximum Likelihood Estimation of Observer Error-Rates.
\newblock \textit{Applied Statistics}, 28(1):20--28, 1979.

\bibitem{patrini2017}
G.\ Patrini et al.
\newblock Making Deep Neural Networks Robust to Label Noise.
\newblock \textit{Proc.\ CVPR}, 2017.

\bibitem{pearl2009}
J.\ Pearl.
\newblock \textit{Causality}, 2nd ed.\ Cambridge University Press, 2009.

\end{thebibliography}

\end{document}
